\documentclass[12pt]{article}

\usepackage{fullpage}
\usepackage{graphicx}
\usepackage{amssymb}
\usepackage{amsmath}
\usepackage[none]{hyphenat}
\usepackage{parskip}
\usepackage[spanish]{babel}
\usepackage[utf8]{inputenc}
\usepackage{hyperref}
\usepackage{fancyhdr}
\usepackage{tasks}
\usepackage{mdframed}
\usepackage{xcolor}
\usepackage{pgfplots}
\usepackage[makeroom]{cancel}
\usepackage{multicol}
\usepackage[shortlabels]{enumitem}
\usepackage{stackrel}
\usepackage{tkz-tab}
\usepackage{xpatch}
\usepackage{tkz-euclide}
\usetkzobj{all}
\xpatchcmd{\tkzTabLine}{$0$}{$\bullet$}{}{}

\setlength{\headheight}{10pt}
\setlength{\headsep}{10pt}
\pagestyle{fancy}
\rhead{\ayudantia \ - \alumno}
\tikzset{t style/.style={style=solid}}

\newcommand*{\mybox}[2]{\colorbox{#1!30}{\parbox{.98\linewidth}{#2}}}

\newenvironment{solucion}
{\begin{mdframed}[backgroundcolor=black!10]
		{\bf Solución:}\\
	}
	{
	\end{mdframed}
}

\newenvironment{alternativas}[1]
{\begin{multicols}{#1}
		\begin{enumerate}[a)]
		}
		{
		\end{enumerate}
	\end{multicols}
}

\newenvironment{preguntas}
{\begin{enumerate}\itemsep12pt
	}
	{
	\end{enumerate}
}

\newcommand{\ayudantia}{{\sc Ayudantía 15}}
\newcommand{\tituloayu}{Compilado Álgebra Lineal}
\newcommand{\fecha}{22 de junio de 2019}
\newcommand{\sigla}{MAT1203}
\newcommand{\nombre}{Álgebra Lineal}
\newcommand{\profesor}{Camilo Perez}
\newcommand{\ano}{2019}
\newcommand{\semestre}{1}
\newcommand{\mail}{mat1203@ifcastaneda.cl}
\newcommand{\alumno}{Ignacio Castañeda - \mail}

\newcommand{\ev}{\Big|}
\newcommand{\ra}{\rightarrow}
\newcommand{\lra}{\leftrightarrow}
\newcommand{\N}{\mathbb{N}}
\newcommand{\R}{\mathbb{R}}
\newcommand{\Exp}[1]{\mathcal{E}_{#1}}
\newcommand{\List}[1]{\mathcal{L}_{#1}}
\newcommand{\EN}{\Exp{\N}}
\newcommand{\LN}{\List{\N}}
\newcommand{\comment}[1]{}
\newcommand{\lb}{\\~\\}
\newcommand{\eop}{_{\square}}
\newcommand{\hsig}{\hat{\sigma}}
\newcommand{\widesim}[2][1.5]{
	\mathrel{\overset{#2}{\scalebox{#1}[1]{$\sim$}}}
}
\newcommand{\wsim}{\widesim{}}
\newcommand{\lh}{\stackrel{L'H}{=}}

\begin{document}
\thispagestyle{empty}

\begin{minipage}{2cm}
	\includegraphics[width=2cm]{../../../../img/logo.pdf}
	\vspace{0.5cm}
\end{minipage}
\begin{minipage}{\linewidth}
	\begin{tabular}{lrl}
		{\scriptsize\sc Pontificia Universidad Catolica de Chile} & \hspace*{0.7in}Curso: &
		\sigla  - \nombre\\
		{\sc Facultad de Matemáticas}&
		Profesor: & \profesor \\
		{\sc Semestre \ano-\semestre} & Ayudante: & {Ignacio Castañeda}\\
		& {Mail:} & \texttt{\mail}
	\end{tabular}
\end{minipage}

\vspace{-10mm}
\begin{center}
	{\LARGE\bf \ayudantia}\\
	\vspace{0.1cm}
	{\tituloayu}\\
	\vspace{0.1cm}
	\fecha\\
	\vspace{0.4cm}
\end{center}

\begin{preguntas}
\item Realiza las siguientes operaciones con los vectores dados
	$$
	v_1 = \begin{pmatrix}
	1\\
	7\\
	8
\end{pmatrix};\qquad
	v_2 = \begin{pmatrix}
	-4\\
	2\\
	-3
\end{pmatrix}; \qquad
	v_3 = \begin{pmatrix}
	0\\
	1\\
	3
\end{pmatrix}; \qquad
	v_4 = \begin{pmatrix}
	13\\
	-3\\
	1
\end{pmatrix},
	 $$
\begin{tasks}(4)
\task $v_1 + v_2$
\task $v_3 - v_4$
\task $v_2 \cdot v_3$
\task $v_1 \times v_2$
\task $v_2 \times v_1$
\task $(v_1 \times v_2) \cdot v_2$
\task $3 - ((2v_1) \cdot v_2)$
\task $v_2 - v_3 \times v_1$
\end{tasks}
\begin{solucion}

\begin{enumerate}[a)]
\item $v_1 + v_2 =\begin{pmatrix}
				1\\
				7\\
				8
				\end{pmatrix} + \begin{pmatrix}
				-4\\
				2\\
				-3
				\end{pmatrix} = \begin{pmatrix}
				-3\\
				9\\
				5
				\end{pmatrix}$
\item $v_3 - v_4 = \begin{pmatrix}
				0\\
				1\\
				3
				\end{pmatrix} - \begin{pmatrix}
				13\\
				-3\\
				1
				\end{pmatrix} =  \begin{pmatrix}
				-13\\
				4\\
				2
				\end{pmatrix}$
\item $v_2 \cdot v_3 = \begin{pmatrix}
				-4\\
				2\\
				-3
				\end{pmatrix} \cdot \begin{pmatrix}
				0\\
				1\\
				3
				\end{pmatrix} = -4 \cdot 0 + 2 \cdot 1 -3 \cdot 3 = -7$
\item $v_1 \times v_2 = \begin{pmatrix}
				1\\
				7\\
				8
			\end{pmatrix} \times \begin{pmatrix}
				-4\\
				2\\
				-3
			\end{pmatrix} = i \left| \begin{matrix} 7 & 8 \\ 2 & -3\end{matrix} \right| - j \left| \begin{matrix} 1 & 8 \\ 4 & -3\end{matrix} \right| + k \left| \begin{matrix} 1 & 7 \\ -4 & 2\end{matrix} \right|$\\
			$$= i (7 \cdot (-3) - 8 \cdot 2) - j (1 \cdot (-3) - 8 \cdot (-4)) + k (1 \cdot 2 - 7 \cdot(-4))$$
			$$= i (-21 -16)) - j(-3--32) + k(2 - - 28)$$
			$$-37 i - 29j + 30k = \begin{pmatrix}
			-37\\
			-29\\
			30
			\end{pmatrix}$$	
\item $v_2 \times v_1 =  \begin{pmatrix}
			-4\\
			2\\
			-3
			\end{pmatrix} \times \begin{pmatrix}
			1\\
			7\\
			8
			\end{pmatrix} = i \left| \begin{matrix} 2 & -3 \\ 7 & 8\end{matrix} \right| - j \left| \begin{matrix} 4 & -3 \\ 1 & 8\end{matrix} \right| + k \left| \begin{matrix} -4 & 2 \\ 1 & 7\end{matrix} \right|$\\
			$$= i (2 \cdot 8 - (-3) \cdot 7) - j ((-4)\cdot 8 - (-3)\cdot 1) + k ((-4) \cdot 7 -2 \cdot 1)$$
			$$= i (16 --21) + j(-32 --3) + k(-28-2)$$
			$$37 i + 29j - 30k = \begin{pmatrix}
			37\\
			29\\
			-30
			\end{pmatrix}$$	
\item $(v_1 \times v_2) \cdot v_2 = \begin{pmatrix}
			-37\\
			-29\\
			30
			\end{pmatrix} \cdot \begin{pmatrix}
			-4\\
			2\\
			-3
			\end{pmatrix} = (-37) \cdot (-4) + (-29) \cdot 2 + 30 \cdot (-3)$
			$$=148 +-58 - 90 = 0$$
\item $3 - ((2v_1) \cdot v_2) = 3 - \left(\left(2\begin{pmatrix}
			1\\
			7\\
			8
			\end{pmatrix}\right) \cdot \begin{pmatrix}
			-4\\
			2\\
			-3
			\end{pmatrix}\right) = 3 - \left(\begin{pmatrix}
			2\\
			14\\
			16
			\end{pmatrix} \cdot \begin{pmatrix}
			-4\\
			2\\
			-3
			\end{pmatrix}\right)$
			$$ = 3 - (2 \cdot (-4) + 14 \cdot 2 + 16 \cdot (-3))$$
			$$3 - (-8 + 28 -48)$$
\item $v_2 - v_3 \times v_1 = \begin{pmatrix}
			-4\\
			2\\
			-3
			\end{pmatrix} -  \begin{pmatrix}
			0\\
			1\\
			3
			\end{pmatrix} \times  \begin{pmatrix}
			1\\
			7\\
			8
			\end{pmatrix}$
			$$= -4i + 2j -3k - (i (8 - 21))- j(0 - 3) + k (0-1))$$
			$$ = -4i + 2j -3k +13i - 3j -k = 9i - j -4k $$
			$$ =\begin{pmatrix}
			9\\
			-4\\
			-2
			\end{pmatrix}$$
\end{enumerate}
\end{solucion}
\item Verificar si los siguientes puntos son o no colineares entre si
\begin{tasks}(2)
\task $P(3,2,5), \ P(0, 1/2, -1), \ P(5, 3, 9)$
\task $P(1,1,4), \ P(-2,-3.-8), \ P(4,4,16)$
\end{tasks}
\begin{solucion}

\begin{enumerate}[a)]
\item $ 
			\begin{pmatrix}
			3\\
			2\\
			5
			\end{pmatrix}; \quad
			\begin{pmatrix}
			0\\
			1/2\\
			-1
			\end{pmatrix}; \quad
			\begin{pmatrix}
			5\\
			3\\
			9
			\end{pmatrix}$\\
			En primer lugar, debemos encontrar dos vectores directores entre dos pares de puntos distintos
			$$d_1 = \begin{pmatrix}
			3\\
			2\\
			5
			\end{pmatrix} - 
			\begin{pmatrix}
			0\\
			1/2\\
			-1
			\end{pmatrix} = \begin{pmatrix}
			3\\
			3/2\\
			6
			\end{pmatrix}$$
			$$d_2 = \begin{pmatrix}
			3\\
			2\\
			5
			\end{pmatrix} - 
			\begin{pmatrix}
			5\\
			3\\
			9
			\end{pmatrix} = \begin{pmatrix}
			-2\\
			-1\\
			-4
			\end{pmatrix}$$
			Dado que es posible representar $d_1$ de la forma $d_1 = \lambda d_2$, con $\lambda = -1,5$, los puntos son colineales.
\item $ 
			\begin{pmatrix}
			1\\
			1\\
			4
			\end{pmatrix}; \quad
			\begin{pmatrix}
			-2\\
			-3\\
			-8
			\end{pmatrix}; \quad
			\begin{pmatrix}
			4\\
			4\\
			16
			\end{pmatrix}$\\
			Igual que antes, buscaremos dos vectores directores
			$$d_1 = \begin{pmatrix}
			1\\
			1\\
			4
			\end{pmatrix} - 
			\begin{pmatrix}
			-2\\
			-3\\
			-8
			\end{pmatrix} = \begin{pmatrix}
			3\\
			4\\
			12
			\end{pmatrix}$$
			$$d_2 = \begin{pmatrix}
			1\\
			1\\
			4
			\end{pmatrix} - 
			\begin{pmatrix}
			4\\
			4\\
			16
			\end{pmatrix} = \begin{pmatrix}
			-3\\
			-3\\
			-12
			\end{pmatrix}$$
			Observando las dos primeras componentes de cada vector, vemos que es imposible representar uno en función del otro, por lo que los puntos no son colineales
\end{enumerate}
\end{solucion}
\item Sean los vectores
		$$
	v_1 = \begin{pmatrix}
	2\\
	1\\
	5
\end{pmatrix};\qquad
	v_2 = \begin{pmatrix}
	-1\\
	-3\\
	0
\end{pmatrix}$$
\begin{enumerate}[a)]
\item Buscar un vector $v_3$ que sea perpendicular a $v_1$ y $v_2$ y verificar que efectivamente sea perpendicular a ambos.
\item Buscar un vector paralelo a $(v_2-v_3)$ y verificar que lo sea
\item Encontrar un vector perpendicular a $(v_3 + 3v_1)$
\end{enumerate}
\begin{solucion}

\begin{enumerate}[a)]
\item Para que sea perpendicular a ambos, basta definir $v_3 = v_1 \times v_2$, es decir
			$$v_3 = \begin{pmatrix}
			2\\
			1\\
			5
			\end{pmatrix} \times \begin{pmatrix}
			-1\\
			-3\\
			0
			\end{pmatrix} = i (0 + 15) - j (0 + 5) + k(-6 + 1) =\begin{pmatrix}
			15\\
			-5\\
			-5
			\end{pmatrix} $$
			Para verificar basta utilizar el producto punto, donde obtendremos que
			$$v_3 \cdot v_1 = 0$$
			$$v_3 \cdot v_2 = 0$$
\item En primer lugar, calculemos el vector $v_2 - v_3$
			$$v_4 = v_2 - v_3 = \begin{pmatrix}
			-1\\
			-3\\
			0
			\end{pmatrix} - \begin{pmatrix}
			15\\
			-5\\
			-5
			\end{pmatrix} = \begin{pmatrix}
			-16\\
			2\\
			5
			\end{pmatrix}$$
			Para obtener un vector paralelo, podemos simplemente ponderar el vector por un escalar, como por ejemplo, 2
			$$v_5 = 2 v_4 = 2 \begin{pmatrix}
			-16\\
			2\\
			5
			\end{pmatrix} = \begin{pmatrix}
			-32\\
			4\\
			10
			\end{pmatrix}$$
			Para verificar utilizaremos el producto cruz, donde obtendremos que
			$$v_5 \times v_4 = 0$$
\item Primero calculemos el vector $v_3 + 3v_1$
			$$v_6 = v_3 + 3v_1 = \begin{pmatrix}
			15\\
			-5\\
			-5
			\end{pmatrix}  + 3 \begin{pmatrix}
			2\\
			1\\
			5
			\end{pmatrix} = \begin{pmatrix}
			21\\
			-2\\
			10
			\end{pmatrix} $$
			Para obtener un vector perpendicular, podemos hacer producto cruz con cualquier otro vector arbitrario. Para hacer los calculos simples, utilizaremos
			$$v_7 = \begin{pmatrix}
			1\\
			0\\
			0
			\end{pmatrix} $$
			Luego, 
			$$v_8 = v_6 \times v_7 = \begin{pmatrix}
			21\\
			-2\\
			10
			\end{pmatrix} \times \begin{pmatrix}
			1\\
			0\\
			0
			\end{pmatrix} = 10j + 2k = \begin{pmatrix}
			0\\
			10\\
			2
			\end{pmatrix}$$
\end{enumerate}
\end{solucion}
\item Detereminar si los siguientes vectores son paralelos, perpendiculares u oblicuos
\begin{tasks}(2)
\task $$ 
			\begin{pmatrix}
			-1\\
			2\\
			5
		\end{pmatrix};\qquad
			\begin{pmatrix}
			2\\
			-4\\
			-10
		\end{pmatrix}$$
\task  $$ 
			\begin{pmatrix}
			2\\
			5\\
			10
		\end{pmatrix};\qquad
			\begin{pmatrix}
			-3\\
			8\\
			-11
		\end{pmatrix}$$
\end{tasks}
\begin{solucion}

\begin{enumerate}[a)]
\item $$ 
			\begin{pmatrix}
			-1\\
			2\\
			5
			\end{pmatrix};\qquad
			\begin{pmatrix}
			2\\
			-4\\
			-10
			\end{pmatrix}$$
			Si se fijan, es evidente que $v_2 = -2v_1$, por lo que comprobaremos de inmediato con el producto cruz para ver si son paralelos
			 $$ 
			\begin{pmatrix}
			-1\\
			2\\
			5
			\end{pmatrix} \times 
			\begin{pmatrix}
			2\\
			-4\\
			-10
			\end{pmatrix} = i (-20 + 20) -j (10 - 10) + k (4 - 4) = \vec{0}$$
			Por lo que ambos vectores son paralelos
\item $$ 
			\begin{pmatrix}
			2\\
			5\\
			10
			\end{pmatrix};\qquad
			\begin{pmatrix}
			-3\\
			8\\
			-11
			\end{pmatrix}$$
			Comenzaremo probando si son perpendiculares con el producto punto
			$$ 
			\begin{pmatrix}
			2\\
			5\\
			10
			\end{pmatrix} \cdot
			\begin{pmatrix}
			-3\\
			8\\
			-11
			\end{pmatrix} = -6 + 40 - 110 = -76$$
			Como es distinto de cero, vemos que no son perpendiculares.\\
			Ahora, utilizaremos el producto cruz para ver si son paralelos
			$$ 
			\begin{pmatrix}
			2\\
			5\\
			10
			\end{pmatrix} \times
			\begin{pmatrix}
			-3\\
			8\\
			-11
			\end{pmatrix} = i(-55 -80) - j (-22 +30) +k(16 + 15) = \begin{pmatrix}
			-135\\
			-8\\
			31
			\end{pmatrix}$$
			Como este no es el vector nulo, tampoco son paralelos. Con esto concluimos que ambos vectores son oblicuos.
\end{enumerate}
\end{solucion}
\item Determinar el plano que pasa por los puntos $P_1(4,-1,-2)$, $P_2(0,0,1)$ y $P_3(2,-3,0)$.
\begin{solucion}
Sabemos que un plano se define de la forma
		$$Ax +By + Cz = D$$
		Reemplazando esto con los tres puntos que tenemos, obtendremos el siguiente sistema de ecuaciones
		$$
		\begin{array}{rcrr}
		4A-B-2C & = & D& \vline\\
		C & = & D & \vline\\
		2A-3B & = & D &\vline\\
		\hline
		\end{array}
		$$
		$$
		\begin{array}{rcrr}
		4A-B-2C & = & C& \vline\\
		2A-3B & = & C &\vline\\
		\hline
		\end{array}
		$$
		$$
		\begin{array}{rcrr}
		4A-B-3C & = & 0& \vline\\
		2A-3B -C& = & 0 &\vline\\
		\hline
		\end{array}
		$$
		$$(1) - 2(2)$$
		$$5B - C = 0 \ra 5B = C$$
		Como hay 4 variables y teníamos 3 ecuaciones, debemos elegir el valor de una de ellas de manera arbitraria.
		$$C = 5$$
		$$\ra D = 5$$
		$$\ra B = 1$$
		$$4A -1 -15 = 0 \ra 4A=16 \ra A=4$$
		Finalmente, el plano es
		$$\Pi:4x + y + 5z = 5$$
\end{solucion}
\item Dado $P_1(0,2,-3)$ y $P_2(1,0,3)$, determinar la recta que pasa por $P_1$ y $P_2$.
\begin{solucion}
Buscamos un vector director
		$$\vec{d} = P_2 - P_1 = (1, -2, 0)$$
		Luego, la recta que pasa por $P_1$ y $P_2$ esta definida por
		$$<x,y,z> = (0, 2, -3) + \lambda (1,-2,0)$$
\end{solucion}
\item Encontrar las ecuaciones de dos planos diferentes cuya intersección sea la recta que pasa por el punto $P_1(1,3,-2)$ y $P_2(2,0,4)$.
\begin{solucion}
Para realizar esto, basta con utilizar dos trios de puntos que contengan a los dos puntos dados y buscar dos planos. La intersección entre estos planos será la recta que pasa por ambos puntos. Para esto, definiremos puntos arbitrarios.\\
		Diremos que
		$$P_3(0, 0, 0), \quad P_4(1,0,0)$$
		Partamos usando los puntos $P_1$, $P_2$ y $P_3$. Sabemos que la ecuación del plano es de la forma
		$$Ax + By + Cz = D$$
		Reemplazando con estos tres puntos, obtenemos el siguiente sistema
		$$
		\begin{array}{rcrr}
		A + 3B - 2C & = & D& \vline\\
		2A + 4C & = & D & \vline\\
		0 & = & D &\vline\\
		\hline
		\end{array}
		$$
		$$
		\begin{array}{rcrr}
		A + 3B - 2C & = & 0& \vline\\
		2A + 4C & = & 0 & \vline\\
		\hline
		\end{array}
		$$
		$$(2) + 2(1)$$
		$$4A + 6B = 0 \ra A = -\dfrac{6B}{4} \ra A = -\dfrac{3B}{2}$$
		Le pondremos un valor arbitrario a $B$,
		$$B = 4$$
		$$\ra A = -6$$
		$$\ra -12 + 4C = 0 \ra C = 3$$
		Por ende,
		$$\Pi_1: -6x + 4y + 3z = 0$$
		Ahora realizamos lo mismo con los puntos $P_1$, $P_2$, $P_4$. Sea la ecuación del plano
		$$Ax + By + Cz = D$$
		Reemplazando con estos tres puntos, obtenemos el siguiente sistema
		$$
		\begin{array}{rcrr}
		A + 3B - 2C & = & D& \vline\\
		2A + 4C & = & D & \vline\\
		A & = & D &\vline\\
		\hline
		\end{array}
		$$
		$$
		\begin{array}{rcrr}
		3B - 2C & = & 0& \vline\\
		A + 4C & = & 0 & \vline\\
		\hline
		\end{array}
		$$
		$$
		\begin{array}{rcrr}
		3B & = & 2C& \vline\\
		A & = & -4C & \vline\\
		\hline
		\end{array}
		$$
		Le daremos un valor arbitrario a $C$
		$$C = 3$$
		$$ \ra B = 2$$
		$$ \ra A = -12 $$
		$$ \ra D = -12 $$
		Con lo que obtenemos el plano
		$$\Pi_2: -12x+2y+3z = -12$$
		Finalmente, los planos buscados son:
		$$\Pi_1: -6x + 4y + 3z = 0$$		
		$$\Pi_2: -12x+2y+3z = -12$$
\end{solucion}
\item Encontrar un plano que sea perpendicular al plano cuya ecuación es $3x -7y +2z = 5$ y que pase por el punto $P(0,2,-1)$
\begin{solucion}
En primer lugar, debemos encontrar el vector normal del plano, que esta dado por los coeficientes de su ecuación, es decir
		$$\vec{n} = \begin{pmatrix}
		3\\-7\\2
		\end{pmatrix}$$
		Si los vectores normales de dos planos son perpendiculares entre si, los planos también lo serán, por lo que basta encontrar un vector perpendicular a $\vec{n}$ y utilizarlo como vector normal para nuestro nuevo plano. Para encontrar un vector perpendicular, realizaremos el producto cruz entre $\vec{n}$ y un vector arbitrario elegido por nosotros
		$$\begin{pmatrix} 3 \\ -7 \\ 2 \end{pmatrix} \times \begin{pmatrix} 1 \\ 1 \\ 1 \end{pmatrix} = \begin{pmatrix} -9 \\ -1 \\ 10 \end{pmatrix}$$
		Para encontrar el plano asociado, basta que utilicemos los coeficientes del vector como coeficientes de la ecuación del plano
		$$-9x -y +10z = D$$
		Por último, para encontrar $D$, reemplazamos el punto que nos dan en la ecuación del plano
		$$-2 -10 = D \ra D = -12$$
		Finalmente,
		$$\Pi: -9x -y +10z = -12$$
\end{solucion}
\item Determinar la ecuación de un plano que pase por $P(4,2,1)$ y que sea paralelo al plano de ecuación $2x-5y+z=6$
\begin{solucion}
De forma análoga, para que dos planos sean paralelos, basta con que sus vectores normales sean paralelos. Dicho de otra forma, lo que va a cambiar será el parametro $D$. Es decir, nuestro plano es
		$$2x - 5y + z = D$$
		Reemplazando con el punto,
		$$8-10+1 = D \ra D = -1$$
		Luego, el plano buscado es
		$$\Pi: 2x - 5y + z = -1$$
\end{solucion}
\item 
\begin{enumerate}[a)]
\item Sean $a=\left(\begin{array}{r}
      1\\-1\\1
    \end{array}\right)$ y $b=\left(\begin{array}{r}
      0\\-1\\2
    \end{array}\right)$, calcule el vector\textbf{ $b-proy_a b$ }y verifique que este es ortogonal a $a$.
\item Demuestre que si $a$ y $b\in\R^3$ el vector $b-proy_a b$ es ortogonal a $a$.
\end{enumerate}
\begin{solucion}

\begin{enumerate}[a)]
\item La proyección de $b$ sobre $a$ corresponde a
$$proy_{a}b = \dfrac{b \cdot a}{|a|^2}a =   \dfrac{\left(\begin{array}{r}
	0\\-1\\2
	\end{array}\right) \cdot\left(\begin{array}{r}
		1\\-1\\1
	\end{array}\right)}{\left(\begin{array}{r}
		1\\-1\\1
	\end{array}\right) \cdot \left(\begin{array}{r}
		1\\-1\\1
	\end{array}\right)}\left(\begin{array}{r}
		1\\-1\\1
	\end{array}\right) = \left(\begin{array}{r}
1\\-1\\1
\end{array}\right)$$
Luego, 
$$b-proy_a b = \left(\begin{array}{r}
0\\-1\\2
\end{array}\right) - \left(\begin{array}{r}
1\\-1\\1
\end{array}\right) = \left(\begin{array}{r}
-1\\0\\1
\end{array}\right)$$
Podemos verificar que el resultado es perpendicular a $a$ utilizando el producto punto, esto es
$$\left(\begin{array}{r}
1\\-1\\1
\end{array}\right) \cdot \left(\begin{array}{r}
-1\\0\\1
\end{array}\right) = 0$$
\item Para demostrar esto, debemos demostrar que
$$(b-proy_a b) \cdot a = 0 $$
Desarrollando,
$$\left(b-\dfrac{b \cdot a}{|a|^2}a\right) \cdot a = 0 $$
$$b \cdot a -\dfrac{b \cdot a}{|a|^2}a \cdot a = 0 $$
$$b \cdot a -\dfrac{b \cdot a}{|a|^2}|a|^2 = 0 $$
$$b \cdot a -b \cdot a = 0 $$
$$0 = 0$$
$$\blacksquare$$
\end{enumerate}
\end{solucion}
\item Determinar si el siguiente sistema es consistente o no
	$$
	\begin{array}{rcr}
	x_1 -6x_2& = & 5\\
	x_2-4x_3+x_4& = & 0\\
	-x_1+6x_2+x_3+5x_4& = & 3\\
	-x_2+5x_3+4x_4 & = & 0
	\end{array}
	$$
\begin{solucion}
En primer lugar, escribiremos el sistema en notación matricial
		$$
		\left[
		\begin{array}{cccc|c}
		1 & -6 & 0 & 0 & 5\\
		0 & 1 & -4 & 1 & 0\\
		-1& 6 & 1 & 5 & 3 \\
		0 & -1 & 5 & 4 &0
		\end{array}
		\right]$$
		A continuación, utilizando operaciones de fila, llevaremos esta matriz a su forma escalonada
		$$\left[
		\begin{array}{cccc|c}
			1 & -6 & 0 & 0 & 5\\
			0 & 1 & -4 & 1 & 0\\
			-1& 6 & 1 & 5 & 3 \\
			0 & -1 & 5 & 4 &0
		\end{array}
		\right] \stackbin[F_4+F_2]{F_3 +F_1}{\wsim}
		\left[
		\begin{array}{cccc|c}
		1 & -6 & 0 & 0 & 5\\
		0 & 1 & -4 & 1 & 0\\
		0 & 0 & 1 & 5 & 8 \\
		0 & 0 & 1 & 5 & 0
		\end{array}
		\right] \stackbin[]{F_4-F_3}{\wsim}
		\left[
		\begin{array}{cccc|c}
		1 & -6 & 0 & 0 & 5\\
		0 & 1 & -4 & 1 & 0\\
		0 & 0 & 1 & 5 & 8 \\
		0 & 0 & 0 & 0 & -8
		\end{array}
		\right] $$
		Aqui podemos apreciar que en el la última fila, todas las variables tienen coeficiente cero y la columna de coeficientes aumentados tiene valor, lo que no es posible, ya que esto significaría que se cumple la igualdad $0 = -8$, lo que no es cierto. Dicho esto, el sistema no es consistente.
\end{solucion}
\item Lleve las siguientes matrices ampliadas a su forma escalonada reducida y determine la existencia y unicidad de las soluciones del sistema.
\begin{tasks}(2)
\task $
		\begin{bmatrix}
		1 & 2 & 4 & 8\\
		2 & 4 & 6 & 8\\
		3 & 6 & 9 & 12
		\end{bmatrix}
		$
\task $
		\begin{bmatrix}
		1 & 2 & 4 & 5\\
		2 & 4 & 5& 4\\
		4 & 5 & 4 & 2
		\end{bmatrix}
		$
\end{tasks}
\begin{solucion}

\begin{enumerate}[a)]
\item $
			\begin{bmatrix}
			1 & 2 & 4 & 8\\
			2 & 4 & 6 & 8\\
			3 & 6 & 9 & 12
			\end{bmatrix}
			\stackbin[F_3 - 3F_1]{F_2-2F_1}{\wsim}
			\begin{bmatrix}
			1 & 2 & 4 & 8\\
			0 & 0 & -2 & -8\\
			0 & 0 & -3 & -12
			\end{bmatrix}
			\stackbin[F_3 \leftarrow -\frac{1}{3}F_3]{F_2 \leftarrow -\frac{1}{2}F_2}{\wsim}
			\begin{bmatrix}
			1 & 2 & 4 & 8\\
			0 & 0 & 1 & 4\\
			0 & 0 & 1 & 4
			\end{bmatrix}$\\\\
			$
			\stackbin[]{F_3- F_2}{\wsim}
			\begin{bmatrix}
			1 & 2 & 4 & 8\\
			0 & 0 & 1 & 4\\
			0 & 0 & 0 & 0
			\end{bmatrix}
			\stackbin[]{F_1- 4F_2}{\wsim}
			\begin{bmatrix}
			1 & 2 & 0 & -8\\
			0 & 0 & 1 & 4\\
			0 & 0 & 0 & 0
			\end{bmatrix}
			$\\\\
			Como hay una variable libre (2 pivotes y 3 filas), existen infinitas soluciones.
\item $
			\begin{bmatrix}
			1 & 2 & 4 & 5\\
			2 & 4 & 5& 4\\
			4 & 5 & 4 & 2
			\end{bmatrix}
			\stackbin[F_3 - 4F_1]{F_2 - 2F_1}{\wsim}
			\begin{bmatrix}
			1 & 2 & 4 & 5\\
			0 & 0 & -3& -6\\
			0 & -3 & -12 & -18
			\end{bmatrix}
			\stackbin[F_3 \leftarrow -\frac{1}{3} F_2]{F_2 \leftarrow -\frac{1}{3}F_3}{\wsim}
			\begin{bmatrix}
			1 & 2 & 4 & 5\\
			0 & 1 & 4& 6\\
			0 & 0 & 1 & 2
			\end{bmatrix}
			$\\\\
			$
			\stackbin[F_2 - 4F_3]{F_1 - 4F_3}{\wsim}
			\begin{bmatrix}
			1 & 2 & 0 & -3\\
			0 & 1 & 0& -2\\
			0 & 0 & 1 & 2
			\end{bmatrix}
			\stackbin[]{F_1 - 2F_2}{\wsim}
			\begin{bmatrix}
			1 & 0 & 0 & 1\\
			0 & 1 & 0& -2\\
			0 & 0 & 1 & 2
			\end{bmatrix}$\\\\
			Como tiene 3 pivotes y 3 filas, el sistema posee solución única
\end{enumerate}
\end{solucion}
\item Sea la matriz $A=
	\begin{bmatrix}
	3 & 5 & -4\\
	-3 & -2 & 4\\
	6 & 1 & -8
	\end{bmatrix}$
\begin{enumerate}[a)]
\item Determinar el conjunto solución de su sistema homogeneo.
\item Describir todas las soluciones de $Ax=b$ con $b=
		\begin{pmatrix}
		7\\
		-1\\
		-4
		\end{pmatrix}$ 
\end{enumerate}
\begin{solucion}

\begin{enumerate}[a)]
\item Determinar el conjunto solución de su sistema homogeneo.\\\\
			El sistema homogeneo corresponde a la ecuación $Ax=\vec{0}$, por lo que nuestra matriz aumentada sería
			$$\left[
			\begin{array}{ccc|c}
			3 & 5 & -4 & 0\\
			-3 & -2 & 4 & 0\\
			6 & 1 & -8 & 0
			\end{array}
			\right] \sim \left[
			\begin{array}{ccc|c}
			3 & 5 & -4 & 0\\
			0 & 3 & 0 & 0\\
			0 & -9 & 0 & 0
			\end{array}
			\right] \sim \left[
			\begin{array}{ccc|c}
			3 & 5 & -4 & 0\\
			0 & 3 & 0 & 0\\
			0 & 0 & 0 & 0
			\end{array}
			\right]$$
			$$\sim \left[
			\begin{array}{ccc|c}
			1 & \frac{5}{3} & -\frac{4}{3} & 0\\
			0 & 1 & 0 & 0\\
			0 & 0 & 0 & 0
			\end{array}
			\right] \sim \left[
			\begin{array}{ccc|c}
			1 & 0 & -\frac{4}{3} & 0\\
			0 & 1 & 0 & 0\\
			0 & 0 & 0 & 0
			\end{array}
			\right]$$
			Esto corresponde al sistema
			$$\begin{array}{rcl}
			x_1 - \frac{4}{3}x_3 & = & 0\\
			x_2 & = & 0\\
			0 & = & 0
			\end{array} \ra \begin{array}{rcl}
			x_1 & = & \frac{4}{3}x_3\\
			x_2 & = & 0\\
			x_3 & = & x_3
			\end{array} \ra x = \begin{pmatrix}
			\frac{4}{3}x_3\\
			0\\
			x_3
			\end{pmatrix} = x_3\begin{pmatrix}
			\frac{4}{3}\\
			0\\
			1
			\end{pmatrix}$$
			Luego, la solución del sistema homogeneo es
			$$S = Gen\left\{\begin{pmatrix}
			\frac{4}{3}\\
			0\\
			1
			\end{pmatrix}\right\}$$
\item Describir todas las soluciones de $Ax=b$ con $b=
			\begin{pmatrix}
			7\\
			-1\\
			-4
			\end{pmatrix}$ \\\\
			Para esto, hacemos lo mismo que en la parte a), pero aumentando la matriz por el vector $b$, es decir
			$$\left[
			\begin{array}{ccc|c}
			3 & 5 & -4 & 7\\
			-3 & -2 & 4 & -1\\
			6 & 1 & -8 & -4
			\end{array}
			\right] \sim \left[
			\begin{array}{ccc|c}
			3 & 5 & -4 & 7\\
			0 & 3 & 0 & 6\\
			0 & -9 & 0 & -18
			\end{array}
			\right] \sim \left[
			\begin{array}{ccc|c}
			3 & 5 & -4 & 7\\
			0 & 3 & 0 & 6\\
			0 & 0 & 0 & 0
			\end{array}
			\right]$$
			$$\sim \left[
			\begin{array}{ccc|c}
			1 & \frac{5}{3} & -\frac{4}{3} & \frac{7}{3}\\
			0 & 1 & 0 & 2\\
			0 & 0 & 0 & 0
			\end{array}
			\right] \sim \left[
			\begin{array}{ccc|c}
			1 & 0 & -\frac{4}{3} & -1\\
			0 & 1 & 0 & 2\\
			0 & 0 & 0 & 0
			\end{array}
			\right]$$
			Esto corresponde al sistema
			$$\begin{array}{rcl}
			x_1 - \frac{4}{3}x_3 & = & -1\\
			x_2 & = & 2\\
			0 & = & 0
			\end{array} \ra \begin{array}{rcl}
			x_1 & = & -1 + \frac{4}{3}x_3\\
			x_2 & = & 2\\
			x_3 & = & x_3
			\end{array}$$
			$$ x = \begin{pmatrix}
			-1 + \frac{4}{3}x_3\\
			2\\
			x_3
			\end{pmatrix} = \begin{pmatrix}
			-1\\2\\0
			\end{pmatrix} + x_3\begin{pmatrix}
			\frac{4}{3}\\
			0\\
			1
			\end{pmatrix}$$
			Luego, la solución del sistema $Ax = b$ es
			$$S = \begin{pmatrix}
			-1\\2\\0
			\end{pmatrix} + Gen\left\{\begin{pmatrix}
			\frac{4}{3}\\
			0\\
			1
			\end{pmatrix}\right\}$$
			Finalmente, podriamos amplificar el vector del generado por 3, para que quede más bonito, con lo que
			$$S = \begin{pmatrix}
			-1\\2\\0
			\end{pmatrix} + Gen\left\{\begin{pmatrix}
			4\\
			0\\
			3
			\end{pmatrix}\right\}$$
\end{enumerate}
\end{solucion}
\item Considere el siguiente sistema de ecuaciones, donde $a$ es una constante.
  $$\begin{array}{llll}
   x_1&+x_2&+x_3&=1 \\
    x_1&+x_2&+ax_3&=1 \\
     ax_1&+ax_2&+x_3&=a\\
      x_1&-ax_2&+ax_3&=0  
  \end{array}$$
\begin{enumerate}[a)]
\item Determine valores de $a$ para los cuales el sistema es inconsistente.
\item Determine valores de $a$ para los cuales el sistema es consistente, y encuentre la solución.
\end{enumerate}
\begin{solucion}
En primer lugar, escribamos el sistema en su forma matricial, esto es
	$$\begin{bmatrix}
	1 & 1 & 1 & 1\\
	1 & 1 & a & 1\\
	a & a & 1 & a\\
	1 &-a & a & 0
	\end{bmatrix}$$
	Ahora, lo que debemos hacer es pivotear la matriz intentando evitar dividir por $a$. En el caso de que sea estrictamente necesario hacerlo, debemos ver el caso donde $a=0$ por separado y luego continuar.
	$$\begin{bmatrix}
	1 & 1 & 1 & 1\\
	1 & 1 & a & 1\\
	a & a & 1 & a\\
	1 &-a & a & 0
	\end{bmatrix} 
	\stackbin[F_4 - F_1]{F_2 - F_1}{\stackbin[]{F_3 - aF_1}{\wsim}}
	\begin{bmatrix}
	1 & 1 & 1 & 1\\
	0 & 0 & a-1 & 0\\
	0 & 0 & 1-a & 0\\
	0 &-1-a & a-1 & -1
	\end{bmatrix} 
	\stackbin[]{F_2 \lra F_4}{\wsim}
	\begin{bmatrix}
	1 & 1 & 1 & 1\\
	0 &-1-a & a-1 & -1\\
	0 & 0 & 1-a & 0\\
	0 & 0 & a-1 & 0
	\end{bmatrix}$$
	$
	\stackbin[]{F_4 + F_3}{\wsim}
	\begin{bmatrix}
	1 & 1 & 1 & 1\\
	0 &-1-a & a-1 & -1\\
	0 & 0 & 1-a & 0\\
	0 & 0 & 0 & 0
	\end{bmatrix}
	\stackbin[]{F_2 + F_3}{\wsim}
	\begin{bmatrix}
	1 & 1 & 1 & 1\\
	0 &-1-a & 0 & -1\\
	0 & 0 & 1-a & 0\\
	0 & 0 & 0 & 0
	\end{bmatrix}$
\begin{enumerate}[a)]
\item 
En la tercera fila no existe ningún valor de $a$ para el cual el sistema es inconsistente, por lo que esta no nos restringe $a$.
Sin embargo, en la segunda fila, si podemos elegir un valor de $a$ para que el sistema sea inconsistente. Esto es,
$$-1-a = 0 \ra a = -1$$
\item Para el resto de valores de $a$, el sistema es consistente, es decir, para
$$a \neq -1$$
Luego, teniendo la matriz en su forma escalonada, obtener la solución del sistema es sencillo.\\
Veamos primero el caso donde $a=1$,
$$\begin{bmatrix}
1 & 1 & 1 & 1\\
0 &-2 & 0 & -1\\
0 & 0 & 0 & 0\\
0 & 0 & 0 & 0
\end{bmatrix} \sim
\begin{bmatrix}
1 & 1 & 1 & 1\\
0 & 1 & 0 & \frac{1}{2}\\
0 & 0 & 0 & 0\\
0 & 0 & 0 & 0
\end{bmatrix} \sim
\begin{bmatrix}
1 & 0 & 1 & \frac{1}{2}\\
0 & 1 & 0 & \frac{1}{2}\\
0 & 0 & 0 & 0\\
0 & 0 & 0 & 0
\end{bmatrix} $$
Luego,
$$\begin{array}{rcl}
x_1 & = & -x_3 + \dfrac{1}{2}\\
x_2 & = & \dfrac{1}{2}\\
x_3 & = & 0 
\end{array} \ra 
x = \begin{pmatrix}
\frac{1}{2} \\
\frac{1}{2} \\
0
\end{pmatrix} + \left<\begin{pmatrix}
-1\\0\\1
\end{pmatrix}\right>$$
Si es que $a \neq 1$, entonces
$$\begin{bmatrix}
1 & 1 & 1 & 1\\
0 &-1-a & 0 & -1\\
0 & 0 & 1-a & 0\\
0 & 0 & 0 & 0
\end{bmatrix}
\stackbin[F_3 \leftarrow \frac{1}{1-a}F_3]{F_2 \leftarrow \frac{1}{-1-a}F_2}{\wsim}
\begin{bmatrix}
1 & 1 & 1 & 1\\
0 & 1 & 0 & \frac{1}{1+a}\\
0 & 0 & 1 & 0\\
0 & 0 & 0 & 0
\end{bmatrix}
\stackbin[F_1-F_3]{F_1 - F_2}{\wsim}
\begin{bmatrix}
1 & 0 & 0 & \frac{a}{1+a}\\
0 & 1 & 0 & \frac{1}{1+a}\\
0 & 0 & 1 & 0\\
0 & 0 & 0 & 0
\end{bmatrix}$$
Luego,
$$\begin{array}{rcl}
x_1 & = & \dfrac{a}{1+a}\\
x_2 & = & \dfrac{1}{1+a}\\
x_3 & = & 0 
\end{array} \ra 
x = \begin{pmatrix}
\frac{a}{1+a} \\
\frac{1}{1+a} \\
0
\end{pmatrix}$$
Notemos que esta última corresponde siempre a una solución única pero que depende de $a$.
\end{enumerate}
\end{solucion}
\item Sea la matriz $A=
	\begin{bmatrix}
	1 & 3 & 4\\
	-4 & 2 & -6\\
	-3 & -2 & -7
	\end{bmatrix}
	$ y $b$ un vector en $R^3$. ¿La ecuación $Ax=b$ es consistente para todo $b$?
\begin{solucion}
Diremos que $b = \begin{pmatrix}
	b_1 \\ b_2 \\ b_3
	\end{pmatrix}$. Luego, el sistema $Ax = b$ se puede representar como
	$$\left[
	\begin{array}{ccc|c}
	1 & 3 & 4 &b_1\\
	-4 & 2 & -6 & b_2\\
	-3 & -2 & -7 & b_3
	\end{array}
	\right] \sim \left[
	\begin{array}{ccc|c}
	1 & 3 & 4 &b_1\\
	0 & 14 & 10 & b_2 + 4b_1\\
	0 & 7 & 5 & b_3 + 3b_1
	\end{array}
	\right]$$
	$$ \sim \left[
	\begin{array}{ccc|c}
	1 & 3 & 4 &b_1\\
	0 & 14 & 10 & b_2 + 4b_1\\
	0 & 0 & 0 & b_3 + 3b_1 - \frac{1}{2}(b_2+4b_1)
	\end{array}
	\right] \sim \left[
	\begin{array}{ccc|c}
	1 & 3 & 4 &b_1\\
	0 & 14 & 10 & b_2 + 4b_1\\
	0 & 0 & 0 & b_1 - \frac{1}{2}b_2 + 4b_3
	\end{array}
	\right] $$
	Luego, el sistema será consistente para $b_1 - \dfrac{1}{2}b_2 + b_3 = 0$. Esto se puede representar como
	$$\begin{array}{rcl}
	b_1 & = & \frac{1}{2}b_2 - b_3\\
	b_2 & = & b_2 \\
	b_3 & = & b_3
	\end{array} \ra b = \begin{pmatrix}
	\frac{1}{2}b_2 - b_3 \\
	b_2\\
	b_3 
	\end{pmatrix} = b_2\begin{pmatrix}
	\frac{1}{2} \\
	1\\
	0 
	\end{pmatrix} +  b_3\begin{pmatrix}
	-1 \\
	0\\
	1 
	\end{pmatrix} $$
	Finalmente,
	$$b = Gen \left\{ \begin{pmatrix}
	\frac{1}{2} \\
	1\\
	0 
	\end{pmatrix}, \begin{pmatrix}
	-1 \\
	0\\
	1
	\end{pmatrix}\right\}$$
\end{solucion}
\item Sean los vectores 
	$$
	v_1 = \begin{pmatrix}
	1\\
	1\\
	2
	\end{pmatrix};\qquad
	v_2 = \begin{pmatrix}
	2\\
	2\\
	-3
	\end{pmatrix}; \qquad
	v_3 = \begin{pmatrix}
	0\\
	1\\
	3
	\end{pmatrix}$$
	determinar si $Gen\{v_1, v_2, v_3\} = R^3$.
\begin{solucion}
Para ver si estos tres vectores generan $\R^3$, tiene que pasar que todos sean L.I. entre si. Una forma de ver esto, es formar una matriz con los vectores y ver la cantidad de pivotes que hay, es decir
		$$\begin{bmatrix}
			\\
			v_1 & v_2 & v_3\\
			&&
		\end{bmatrix} \sim 
		\begin{bmatrix}
		1 & 2 & 0\\
		1 & 2 & 1\\
		2 & -3 & 3
		\end{bmatrix} \sim
		\begin{bmatrix}
		1 & 2 & 0\\
		0 & 0 & 1\\
		0 & -7 & 3
		\end{bmatrix} \sim
		\begin{bmatrix}
		1 & 2 & 0\\
		0 & -7 & 3\\
		0 & 0 & 1
		\end{bmatrix}$$
		Como la matriz tiene 3 pivotes, quiere decir que hay 3 vectores L.I., es decir, todos son linealmente independientes entre si, formando asi $\R^3$
\end{solucion}
\item Sean 
$$v_1=\left(\begin{array}{r}
  1\\0\\-2
\end{array}\right), \quad v_2=\left(\begin{array}{r}
  -2\\1\\7
\end{array}\right) \quad y \quad v_3=\left(\begin{array}{r}
  h\\0\\-2
\end{array}\right)$$
¿Para qué valor(es) de $h$ $Gen\{v_1, v_2, v_3\}=Gen\{v_1,v_2\}$ ?
\begin{solucion}
Para que $Gen\{v_1, v_2, v_3\}=Gen\{v_1,v_2\}$, debe ocurrir que estos generados tengan la misma cantidad de vectores L.I. o dicho de otra forma, al poner estos vectores en una matriz y escalonarla, las matrices deben tener la misma cantidad de pivotes.\\
	\\
	Comencemos por $Gen\{v_1,v_2\}$, dado que $h$ no influye en este generado.
	$$\begin{bmatrix}
	\\
	v_1 & v_2\\
	&&
	\end{bmatrix} \sim 
	\begin{bmatrix}
	1 & -2\\
	0 & 1\\
	-2 & 7
	\end{bmatrix} \sim 
	\begin{bmatrix}
	1 & -2\\
	0 & 1\\
	0 & 3
	\end{bmatrix} \sim 
	\begin{bmatrix}
	1 & -2\\
	0 & 1\\
	0 & 0
	\end{bmatrix}$$
	Notemos que esta matriz tiene 2 pivotes, por lo que debemos elegir $h$, de tal manera que $Gen\{v_1, v_2, v_3\}$ también tenga 2 pivotes.\\
	\\
	Procedamos ahora con $Gen\{v_1, v_2, v_3\}$,
	$$\begin{bmatrix}
	\\
	v_1 & v_2 & v_3\\
	&&
	\end{bmatrix} \sim 
	\begin{bmatrix}
	1 & -2 & h\\
	0 & 1 & 0\\
	-2 & 7 & -2
	\end{bmatrix} \sim 
	\begin{bmatrix}
	1 & -2 & h\\
	0 & 1 & 0\\
	0 & 3 & 2h-2
	\end{bmatrix} \sim 
	\begin{bmatrix}
	1 & -2 & h\\
	0 & 1 & 0\\
	0 & 0 & 2h-2
	\end{bmatrix}$$
	Luego, para que la matriz tenga 2 pivotes, debe ocurrir que
	$$2h-2 = 0 \ra h = 1$$
\end{solucion}
\item Sean $\{u, v, w\}$ un conjunto de vectores linealmente independientes. Demuestre que el conjunto $\{u+v, u+2w, v+3u+w\}$ es linealmente independiente.
\begin{solucion}
$$P.D. \quad \{y, v, w\}\quad L.I. \ra \{u+v, u+2w, v+3u+w\}\quad L.I.$$
		Para que $\{u+v, u+2w, v+3u+w\}$ sea L.I., tiene que cumplirse que el sistema
		$$\alpha(u+v) + \beta(u+2w) + \gamma(v+3u+w) = 0$$
		tenga solucion única $\begin{pmatrix}
		\alpha\\ \beta \\ \gamma
		\end{pmatrix} = \begin{pmatrix} 0\\0\\0\end{pmatrix}$\\
		Trabajemos entonces con el sistema
		$$\alpha(u+v) + \beta(u+2w) + \gamma(v+3u+w) = 0$$
		$$\alpha u+ \alpha v + \beta u+2\beta w + \gamma v+3\gamma u+\gamma w = 0$$
		$$(\alpha + \beta + 3 \gamma)u + (\alpha + \gamma)v + (2\beta + \gamma)w = 0$$
		Recordemos que ${u, v, w}$ es L.I, por lo que debe cumplirse que
		$$\begin{array}{rcl}
		\alpha + \beta + 3 \gamma & = & 0\\
		\alpha + \gamma & = & 0\\
		2\beta + \gamma & = & 0
		\end{array}$$
		Este sistema lo podemos expresar de forma matricial,
		$$\begin{bmatrix}
		1 & 1 & 3\\
		1 & 0 & 1 \\
		0 & 2 & 1
		\end{bmatrix} \stackrel{F.E.}{\sim} \begin{bmatrix}
		1 & 0 & 1\\
		0 & 1 & -2\\
		0 & 0 & 1
		\end{bmatrix}$$
		Cuya solución es
		$$\begin{array}{rcl}
		\alpha & = & 0\\
		\beta & = & 0\\
		\gamma & = & 0
		\end{array}$$
		$$q.e.d$$
	
\end{solucion}
\item Demuestre que el conjunto $\{u, v, w\}$ es L.I. si y solo si el conjunto $\{u+v, u+w, v+w\}$ es L.I.
\begin{solucion}
Como nos dicen si y solo si, debemos demostrar en ambas direcciones
		$$(\Longrightarrow)$$
		Digamos que $\{u, v, w\}$ es L.I.\\
		\\
		P.D. $\{u+v, u+w, v+w\}$ es L.I.\\
		\\
		Para que $\{u+v, u+w, v+w\}$ sea L.I., el sistema
		$$\alpha (u+v) + \beta (u+w) + \gamma (v+w) = 0$$
		debe tener solución única $\alpha = 0$, $\beta = 0$, $\gamma = 0$
		Reordenando,
		$$\alpha (u+v) + \beta (u+w) + \gamma (v+w) = 0$$
		$$(\alpha + \beta) u + (\alpha + \gamma) v + (\beta + \gamma) w = 0$$
		Como sabemos que el conjunto $\{u, v, w\}$ es L.I., entonces tenemos que
		$$\begin{array}{rl}
		\alpha + \beta & = 0\\
		\alpha + \gamma & = 0\\
		\beta + \gamma & = 0
		\end{array}$$
		Este sistema tiene por solución
		$$\alpha = 0, \quad \beta = 0, \quad \gamma = 0$$
		Con lo que demostramos que $\{u+v, u+w, v+w\}$ es L.I.
		$$(\Longleftarrow)$$
		Digamos que $\{u+v, u+w, v+w\}$ es L.I.\\
		\\
		P.D. $\{u, v, w\}$ es L.I.\\
		\\
		Para que $\{u, v, w\}$ sea L.I., el sistema
		$$c_1u + c_2v + c_3w$$
		debe tener solución única $c_1 = 0$, $c_2 = 0$, $c_3 = 0$\\
		\\
		Como $\{u+v, u+w, v+w\}$ es L.I., sabemos que el sistema
		$$\alpha (u+v) + \beta (u+w) + \gamma (v+w) = 0$$
		tiene solución única $\alpha = 0$, $\beta = 0$, $\gamma = 0$
		Reordenando,
		$$\alpha (u+v) + \beta (u+w) + \gamma (v+w) = 0$$
		$$(\alpha + \beta) u + (\alpha + \gamma) v + (\beta + \gamma) w = 0$$
		Con lo que tenemos que
		$$c_1 = \alpha + \beta = 0, \quad c_2 = \alpha + \gamma = 0, \quad c_3 = \beta + \gamma = 0$$
		es la única solución del sistema. \\
		\\
		Luego, $\{u, v, w\}$ es L.I.
		$$q.e.d$$
\end{solucion}
\item Sea $L: P_2(\R) \ra \R^2$ una transformación lineal tal que
	$$L(1+x) = \begin{pmatrix}
	1\\
	1
	\end{pmatrix}, \quad L(1+x+x^2) = \begin{pmatrix}
	1\\
	-1
	\end{pmatrix}\ y \ L(1+2x) = \begin{pmatrix}
	1\\
	2
	\end{pmatrix}$$
	determine $L(a+bx+cx^2)$ para todo $a,b,c \in \R$.
\begin{solucion}
Recordemos la siguientes propiedades de las transformaciones lineales
		$$L(v_1) + L(v_2) = L(v_1 + v_2), \quad \alpha L(v) = L(\alpha v)$$
		Debemos buscar 
		$$L(1), \quad L(x), \quad L(x^2)$$
		Para esto, debemos jugar con la información que nos dan hasta encontrar cada uno de ellos.
		\begin{center}\rule{14.5cm}{0.1pt}\end{center}
		$$L(1+x+x^2) - L(1+x) = \begin{pmatrix}
		1\\-1
		\end{pmatrix} - \begin{pmatrix}
		1\\1
		\end{pmatrix}$$
		$$L(x^2) = \begin{pmatrix}
		0\\-2
		\end{pmatrix}$$
		
		\begin{center}\rule{14.5cm}{0.1pt}\end{center}
		$$L(1+2x) - L(1+x) = \begin{pmatrix}
		1\\2
		\end{pmatrix} - \begin{pmatrix}
		1\\1
		\end{pmatrix}$$
		$$L(x) = \begin{pmatrix}
		0\\1
		\end{pmatrix}$$
		
		\begin{center}\rule{14.5cm}{0.1pt}\end{center}
		$$L(1+x) - L(x) = \begin{pmatrix}
		1\\1
		\end{pmatrix} - \begin{pmatrix}
		0\\1
		\end{pmatrix}$$
		$$L(x) = \begin{pmatrix}
		1\\0
		\end{pmatrix}$$
		
		Luego, 
		$$L(a+bx+cx^2) = aL(1) + bL(x) + cL(x^2)$$
		$$L(a+bx+cx^2) = aL\begin{pmatrix}1\\0\end{pmatrix} + b\begin{pmatrix}0\\1\end{pmatrix} + c\begin{pmatrix}0\\-2\end{pmatrix}$$
		$$L(a+bx+cx^2) = \begin{pmatrix}a\\b-2c\end{pmatrix}$$
\end{solucion}
\item Sea $T$ una transformación lineal tal que
	$$T\begin{bmatrix}2 \\ 1\end{bmatrix} = \begin{bmatrix}1 \\ 0 \\ 1 \\ 0\end{bmatrix} 
	\quad y \quad
	T \begin{bmatrix}3 \\ 1\end{bmatrix}  =\begin{bmatrix}0 \\ 1 \\ 0 \\ 1\end{bmatrix}$$
	Determine la matriz que representa a $T$
\begin{solucion}
Estamos buscando la matriz $\begin{bmatrix}
		T\begin{bmatrix}
		1 \\ 0
		\end{bmatrix}
		&
		
		T\begin{bmatrix}
		0 \\ 1
		\end{bmatrix}
		\end{bmatrix}$.\\
		Entonces, lo que necesitamos es buscar $T\begin{bmatrix}
		1 \\ 0
		\end{bmatrix}$ y $T\begin{bmatrix}
			0 \\ 1
		\end{bmatrix}$ y para ello, debemos escribir cada uno como una combinación lineal de $T\begin{bmatrix}
		2 \\ 1
		\end{bmatrix}$ y $T\begin{bmatrix}
		3 \\ 1
		\end{bmatrix}$, es decir
		$$\begin{bmatrix}
		1 \\ 0
		\end{bmatrix} = \alpha \begin{bmatrix}
		2 \\ 1
		\end{bmatrix} +  \beta \begin{bmatrix}
		3 \\ 1
		\end{bmatrix} \quad y \quad
		\begin{bmatrix}
		0 \\ 1
		\end{bmatrix} = \gamma \begin{bmatrix}
		2 \\ 1
		\end{bmatrix} +  \delta \begin{bmatrix}
		3 \\ 1
		\end{bmatrix}$$
		O sea, debemos resolver los sistemas
		$$\begin{array}{rcl}
		2\alpha + 3\beta & = & 1\\
		\alpha + \beta & = & 0
		\end{array} \qquad \qquad \qquad
		\begin{array}{rcl}
		2\gamma + 3\delta & = & 0\\
		\gamma + \delta & = & 1
		\end{array}$$
		La solución de estos sistemas es
		$$\alpha = -1, \quad \beta = 1, \quad \gamma = 3, \quad \delta = -2$$
		Por lo que
		$$\begin{bmatrix}
		1 \\ 0
		\end{bmatrix} = \begin{bmatrix}
		3 \\ 1
		\end{bmatrix}- \begin{bmatrix}
		2 \\ 1
		\end{bmatrix}  \quad y \quad
		\begin{bmatrix}
		0 \\ 1
		\end{bmatrix} = 3 \begin{bmatrix}
		2 \\ 1
		\end{bmatrix} - 2 \begin{bmatrix}
		3 \\ 1
		\end{bmatrix}$$
		Luego,
		$$T\begin{bmatrix}
		1 \\ 0
		\end{bmatrix} = T\begin{bmatrix}
		3 \\ 1
		\end{bmatrix}- T\begin{bmatrix}
		2 \\ 1
		\end{bmatrix} = \begin{bmatrix}0 \\ 1 \\ 0 \\ 1\end{bmatrix} - \begin{bmatrix}1 \\0 \\ 1 \\ 0\end{bmatrix} = \begin{bmatrix}-1 \\ 1 \\ -1 \\ 1\end{bmatrix}$$
		y
		$$T\begin{bmatrix}
		0 \\ 1
		\end{bmatrix} = 3 T\begin{bmatrix}
		2 \\ 1
		\end{bmatrix} - 2 T\begin{bmatrix}
		3 \\ 1
		\end{bmatrix} = 3 \begin{bmatrix}1 \\ 0 \\ 1 \\ 0\end{bmatrix} - 2 \begin{bmatrix}0 \\ 1 \\ 0 \\ 1\end{bmatrix} = \begin{bmatrix}3 \\ -2 \\ 3 \\ -2\end{bmatrix}$$
		Finalmente, la matriz que representa $T$ es
		$$\begin{bmatrix}
		T\begin{bmatrix}
		1 \\ 0
		\end{bmatrix}
		&
		
		T\begin{bmatrix}
		0 \\ 1
		\end{bmatrix}
		\end{bmatrix} =
		\begin{bmatrix}
		-1 & 3 \\
		1 & -2 \\
		-1 & 3 \\
		1 & -2
		\end{bmatrix}$$
\end{solucion}
\item Sea la transformación lineal $T$, dada por la matriz
	$$T = \begin{bmatrix}
	1 & 2 & -1\\
	2 & 4 & -2
	\end{bmatrix}$$
	Determine la imagen del plano de ecuación $x_1 + x_3 = 1$ por T.
\begin{solucion}
En primer lugar, escribimos el plano en su forma paramétrica
		$$\begin{array}{rcl}
		x_1 & = &1-x_3\\
		x_2 & = &x_2\\
		x_3 & = &x_3
		\end{array}$$
		Vemos entonces que 3 puntos arbitrarios pertenecienctes al plano son
		$$P_1(1,0,0), \quad P_2(1,1,0), \quad P_3(0,0,1) $$
		Luego, dos vectores directores que definen al plano son
		$$\vec{d_1} = \vec{P_1P_2} = (0,1,0), \quad =\vec{d_2} = \vec{P_1P_3} = (-1,0,1)$$
		Con esto, podemos expresar el plano en su forma vectorial
		$$\vec{P}= \vec{P_1} + \alpha \vec{d_1}  + \beta \vec{d_2}$$
		Para encontrar la imagen del plano por $T$, basta con hacer
		$$T\vec{P}= T\vec{P_1} + \alpha T\vec{d_1}  + \beta T\vec{d_2}$$
		Como
		$$T\vec{P_1} = \begin{bmatrix}
		1 & 2 & -1\\
		2 & 4 & -2
		\end{bmatrix} \begin{bmatrix}
		1 \\ 0 \\ 0
		\end{bmatrix} = \begin{bmatrix}
		1 \\ 2
		\end{bmatrix}$$
		$$T\vec{d_1} = \begin{bmatrix}
		1 & 2 & -1\\
		2 & 4 & -2
		\end{bmatrix} \begin{bmatrix}
		0 \\ 1 \\ 0
		\end{bmatrix} = \begin{bmatrix}
		2 \\ 4
		\end{bmatrix}$$
		$$T\vec{d_2} = \begin{bmatrix}
		1 & 2 & -1\\
		2 & 4 & -2
		\end{bmatrix} \begin{bmatrix}
		-1 \\ 0 \\ 1
		\end{bmatrix} = \begin{bmatrix}
		-2 \\ -4
		\end{bmatrix}$$
		Tenemos que
		$$TP =  \begin{bmatrix}
		1 \\ 2
		\end{bmatrix} + \alpha  \begin{bmatrix}
		2 \\ 4
		\end{bmatrix} + 
		\beta  \begin{bmatrix}
		-2 \\ -4
		\end{bmatrix} =  \begin{bmatrix}
		1 + 2\alpha  - 2\beta \\
		2 + 4\alpha - 5\beta
		\end{bmatrix} = 
		(1+2\alpha - 2\beta)  \begin{bmatrix}
		1 \\ 2
		\end{bmatrix}$$
		En otras palabras, la imagen del plano por la traslación $T$ corresponde a los ponderados de $\begin{bmatrix}
			1 \\ 2
		\end{bmatrix}$, es decir
		$$T\vec{P} = Gen\left\{ \begin{bmatrix}
		1 \\ 2
		\end{bmatrix} \right\}$$
\end{solucion}
\item Sea $\mathbb{P} \subset \R^3$ el plano de ecuación $x_1-x_2+1=0$ y sea $T:\R^3 \ra \R^4$ la transformación dada por $T(x) = Ax$, donde
$$A = \begin{bmatrix}
1 & 0 & 3\\
-1 & 1 & 0 \\
0 & 1 & 3\\
0 & -2 & -6
\end{bmatrix}$$
Demuestre que la imagen del plano $\mathbb{P}$ bajo la transformación $T$ es una recta en $\R^4$.
\begin{solucion}
Despejando $x_2$ en la ecuación del plano, obtenemos $x_2 = x_1 + 1$.\\

Reemplazando en la transformación lineal,
{\small $$T(x) = T(x_1, x_2, x_3) = T(x_1, x_1+1, x_3) = \begin{bmatrix}
1 & 0 & 3\\
-1 & 1 & 0 \\
0 & 1 & 3\\
0 & -2 & -6
\end{bmatrix}
\begin{pmatrix}
x_1 \\
x_1+1\\
x_3
\end{pmatrix} = 
\begin{pmatrix}
x_1 +3x_3\\
1\\
1+x_1+3x_3\\
-2-2x_1-6x_3
\end{pmatrix}$$}
Esto lo podemos descomponer en
$$T(x) = 
\begin{pmatrix}
x_1 +3x_3\\
1\\
1+x_1+3x_3\\
-2-2x_1-6x_3
\end{pmatrix} =
\begin{pmatrix}
0\\
1\\
1\\
-2
\end{pmatrix} +
x_1 \begin{pmatrix}
1\\
0\\
1\\
-2
\end{pmatrix} +
x_3 \begin{pmatrix}
3\\
0\\
3\\
-6
\end{pmatrix}
$$
Notemos ahora que el vector asociado a $x_3$ es el triple del vector asociado a $x_1$, por lo que podemos factorizarlo de la siguiente forma
$$T(x) = 
\begin{pmatrix}
0\\
1\\
1\\
-2
\end{pmatrix} +
x_1 \begin{pmatrix}
1\\
0\\
1\\
-2
\end{pmatrix} +
3x_3 \begin{pmatrix}
1\\
0\\
1\\
-2
\end{pmatrix} = 
\begin{pmatrix}
0\\
1\\
1\\
-2
\end{pmatrix} +
(x_1 + 3x_3) \begin{pmatrix}
1\\
0\\
1\\
-2
\end{pmatrix}
$$
{\small Esto corresponde a la ecuación de una recta con dirección $\begin{pmatrix}
1\\
0\\
1\\
-2
\end{pmatrix}$ y que pasa por $\begin{pmatrix}
0\\
1\\
1\\
-2
\end{pmatrix}$.}
\end{solucion}
\item Determine si las siguientes afirmaciones son verdaderas o falsas.
\begin{enumerate}[a)]
\item Para todo $u, v \in \R^n$ se tiene que $||u+v|| < ||u||+||v||$
\item Si $u \in Gen\{v_1, v_2\}$ entonces $Gen\{u, v_1\} \subset Gen\{u, v_2\}$
\item Sea $A$ una matriz de $n \times m$ tal que $n < m $, entonces las columnas de A son L.D.
\item Si $A$ es una matriz de $2 \times 3$ tal que $Ax = \begin{pmatrix}
		1\\
		1
		\end{pmatrix}$ y $ Ax = \begin{pmatrix}
		1\\
		2
		\end{pmatrix}$ tienen solución, entonces las filas de $A$ son L.I.
\end{enumerate}
\begin{solucion}

\begin{enumerate}[a)]
\item Para todo $u, v \in \R^n$ se tiene que $||u+v|| < ||u||+||v||$\\\\
			Tomemos $u = v = \vec{0}$. Luego,
			$$||\vec{0} + \vec{0}|| < ||\vec{0}|| + ||\vec{0}||$$
			$$0 < 0$$
			Lo que es falso. Finalmente, la afirmación es FALSA
\item Si $u \in Gen\{v_1, v_2\}$ entonces $Gen\{u, v_1\} \subset Gen\{u, v_2\}$\\\\
			Tomemos
			$$u = \begin{pmatrix}
			0 \\ 0 
			\end{pmatrix}, \quad v_1 = \begin{pmatrix}
			1 \\ 0 
			\end{pmatrix}, \quad v_2 = \begin{pmatrix}
			0 \\ 1
			\end{pmatrix}$$
			Aqui se cumple que $u \in Gen\{v_1, v_2\}$. Sin embargo,
			$$Gen\left\{\begin{pmatrix}
			0 \\ 0 
			\end{pmatrix}, \begin{pmatrix}
			1 \\ 0 
			\end{pmatrix}\right\} \subset Gen\left\{\begin{pmatrix}
			0 \\ 0 
			\end{pmatrix}, \begin{pmatrix}
			0 \\ 1
			\end{pmatrix}\right\}$$
			$$Gen\left\{\begin{pmatrix}
			1 \\ 0 
			\end{pmatrix}\right\} \subset Gen\left\{\begin{pmatrix}
			0 \\ 1
			\end{pmatrix}\right\}$$
			Lo que claramente no es cierto. Luego, la afirmación es FALSA
\item Sea $A$ una matriz de $n \times m$ tal que $n < m $, entonces las columnas de A son L.D.\\\\
			Esto significa que la matriz esta compuesta por $m$ vectores en $\R^n$. Dicho esto, notemos que pueden haber a lo mas $n$ vectores L.I. (máximo $n$ pivotes) Como $m>n$, deben haber vectores L.D., por lo que la afirmación es VERDADERA.
\item Si $A$ es una matriz de $2 \times 3$ tal que $Ax = \begin{pmatrix}
			1\\
			1
			\end{pmatrix}$ y $ Ax = \begin{pmatrix}
			1\\
			2
			\end{pmatrix}$ tienen solución, entonces las filas de $A$ son L.I.\\\\
			Asumiendo que ambos sistemas tienen solución, digamos que las filas de $A$ no serán L.I., es decir, una será múltiplo de la otra. Esto se puede expresar como
			$$A = \begin{bmatrix}
			a_{11} & a_{12} & a_{13} \\
			a_{21} & a_{22} & a_{23}
			\end{bmatrix} = \begin{bmatrix}
			a_{11} & a_{12} & a_{13} \\
			\lambda a_{11} & \lambda a_{12} & \lambda a_{13}
			\end{bmatrix}$$
			Resolvamos ahora el sistema $Ax = \begin{pmatrix}
			b_1 \\ b_2
			\end{pmatrix}$
			$$\begin{array}{lcl}
			a_{11}x_1 + a_{12}x_2 + a_{13}x_3 & = & b_1 \\
			a_{21}x_1 + a_{22}x_2 + a_{23}x_3 & = & b_2			
			\end{array} \ra \begin{array}{lcl}
			a_{11}x_1 + a_{12}x_2 + a_{13}x_3 & = & b_1 \\
			\lambda a_{11}x_1 + \lambda a_{12}x_2 + \lambda a_{13}x_3 & = & \lambda b_1			
			\end{array}$$
			Aqui podemos ver que a lo más uno de los dos sistemas puede tener solución, lo que contradice nuestra hipotesis inicial. Por esto, las filas de $A$ deben ser L.I. Con esto concluimos que la afirmación es VERDADERA
\end{enumerate}
\end{solucion}
\item Sea
	$$A = \begin{bmatrix}
	1 & 0 & 6\\
	-1 & 2 & 0\\
	0 & 5 & -1
	\end{bmatrix}$$
	Escriba $A$ como un producto de matrices elementales
\begin{solucion}
Para esto, escalonamos la matriz $A$ hasta reducirla a la matriz identidad. Es importante que hagamos esto paso por paso, sin realizar mas de una operación por fila a la vez.
		$$\begin{bmatrix}
		1 & 0 & 6\\
		-1 & 2 & 0\\
		0 & 5 & -1
		\end{bmatrix} \widesim{F_2+F_1}
		\begin{bmatrix}
		1 & 0 & 6\\
		0 & 2 & 6\\
		0 & 5 & -1
		\end{bmatrix} \widesim{F_3-\frac{5}{2}F_2}
		\begin{bmatrix}
		1 & 0 & 6\\
		0 & 2 & 6\\
		0 & 0 & -16
		\end{bmatrix} \widesim{F_3 \leftarrow -\frac{1}{16}F_3}
		\begin{bmatrix}
		1 & 0 & 6\\
		0 & 2 & 6\\
		0 & 0 & 1
		\end{bmatrix}$$
		$$\widesim{F_2 \leftarrow \frac{1}{2}F_2}
		\begin{bmatrix}
		1 & 0 & 6\\
		0 & 1 & 3\\
		0 & 0 & 1
		\end{bmatrix} \widesim{F_2 - 3F_3}
		\begin{bmatrix}
		1 & 0 & 6\\
		0 & 1 & 0\\
		0 & 0 & 1
		\end{bmatrix} \widesim{F_1 - 6F_3}
		\begin{bmatrix}
		1 & 0 & 0\\
		0 & 1 & 0\\
		0 & 0 & 1
		\end{bmatrix}$$
		Notemos que todas estas operaciones las podemos escribir como matrices elementales, aplicando la operación a la matriz identidad Sean
		$$E_1 = \begin{bmatrix}
		1 & 0 & 0\\
		1 & 1 & 0\\
		0 & 0 & 1
		\end{bmatrix}, \quad E_2 = \begin{bmatrix}
		1 & 0 & 0\\
		0 & 1 & 0\\
		0 & -\frac{5}{2} & 1
		\end{bmatrix}, E_3 = \begin{bmatrix}
		1 & 0 & 0\\
		0 & 1 & 0\\
		0 & 0 & -\frac{1}{16}
		\end{bmatrix}$$
		$$E_4 = \begin{bmatrix}
		1 & 0 & 0\\
		0 & \frac{1}{2} & 0\\
		0 & 0 & 1
		\end{bmatrix}, \quad E_5 = \begin{bmatrix}
		1 & 0 & 0\\
		0 & 1 & -3\\
		0 & 0 & 1
		\end{bmatrix}, E_6 = \begin{bmatrix}
		1 & 0 & -6\\
		0 & 1 & 0\\
		0 & 0 & 1
		\end{bmatrix}$$
		tenemos que
		$$E_6E_5E_4E_3E_2E_1A = I$$
		por lo que
		$$A = E_1^{-1}E_2^{-1}E_3^{-1}E_4^{-1}E_5^{-1}E_6^{-1}$$
		donde
		$$E_1^{-1} = \begin{bmatrix}
		1 & 0 & 0\\
		-1 & 1 & 0\\
		0 & 0 & 1
		\end{bmatrix}, \quad E_2^{-1} = \begin{bmatrix}
		1 & 0 & 0\\
		0 & 1 & 0\\
		0 & \frac{5}{2} & 1
		\end{bmatrix}, E_3^{-1} = \begin{bmatrix}
		1 & 0 & 0\\
		0 & 1 & 0\\
		0 & 0 & -16
		\end{bmatrix}$$
		$$E_4^{-1} = \begin{bmatrix}
		1 & 0 & 0\\
		0 & 2 & 0\\
		0 & 0 & 1
		\end{bmatrix}, \quad E_5^{-1} = \begin{bmatrix}
		1 & 0 & 0\\
		0 & 1 & 3\\
		0 & 0 & 1
		\end{bmatrix}, E_6^{-1} = \begin{bmatrix}
		1 & 0 & 6\\
		0 & 1 & 0\\
		0 & 0 & 1
		\end{bmatrix}$$
		Finalmente,
		$$A = \begin{bmatrix}
		1 & 0 & 0\\
		-1 & 1 & 0\\
		0 & 0 & 1
		\end{bmatrix}\begin{bmatrix}
		1 & 0 & 0\\
		0 & 1 & 0\\
		0 & \frac{5}{2} & 1
		\end{bmatrix}\begin{bmatrix}
		1 & 0 & 0\\
		0 & 1 & 0\\
		0 & 0 & \frac{1}{16}
		\end{bmatrix}\begin{bmatrix}
		1 & 0 & 0\\
		0 & -\frac{1}{2} & 0\\
		0 & 0 & 1
		\end{bmatrix}\begin{bmatrix}
		1 & 0 & 0\\
		0 & 1 & 3\\
		0 & 0 & 1
		\end{bmatrix}\begin{bmatrix}
		1 & 0 & 6\\
		0 & 1 & 0\\
		0 & 0 & 1
		\end{bmatrix}$$
\end{solucion}
\item Determinar la matriz inversa de
$$A = \begin{bmatrix}
1 & -1 & 0 \\
0 & 1 & 0 \\
2 & 0 & 1
\end{bmatrix}$$
\begin{solucion}
Para obtener $A^{-1}$, ampliamos con la matriz identidad y pivoteamos hasta obtener la matriz identidad al lado izquierdo. Luego, la matriz resultante en el lado derecho será la inversa.\\
\\
Procedamos,
$$\left[
\begin{array}{ccc|ccc}
1 &-1 & 0 & 1 & 0 & 0 \\
0 & 1 & 0 & 0 & 1 & 0 \\
2 & 0 & 1 & 0 & 0 & 1
\end{array}
\right] \sim
\left[
\begin{array}{ccc|ccc}
1 &-1 & 0 & 1 & 0 & 0 \\
0 & 1 & 0 & 0 & 1 & 0 \\
0 & 0 & 1 &-2 & 0 & 1
\end{array}
\right] \sim
\left[
\begin{array}{ccc|ccc}
1 & 0 & 0 & 1 &-1 & 0 \\
0 & 1 & 0 & 0 & 1 & 0 \\
0 & 0 & 1 &-2 & 0 & 1
\end{array}
\right]
$$
Finalmente,
$$A^{-1} = \begin{bmatrix}
 1 &-1 & 0 \\
 0 & 1 & 0 \\
-2 & 0 & 1
\end{bmatrix}$$
\end{solucion}
\item Sea $A$ una matriz de $n\times m$ con columnas LI. Demuestre que si $P=A(A^tA)^{-1}A^{t}$, entonces $P^2=P$ y $P=P^t$.
\begin{solucion}
	Recordemos que
	$$(AB)^t = B^tA^t \qquad y \qquad (AB)^{-1} = B^{-1}A^{-1}$$
	En primer lugar, demostremos que 
	$$P^2 = P \ra P = P^2$$
	Reemplazamos con $P=A(A^tA)^{-1}A^{t}$, es decir
	$$\begin{array}{rcl}
	P & = & P^2 \\
	  & = & (A(A^tA)^{-1}A^{t})(A(A^tA)^{-1}A^{t}) \\
	  & = & A(A^tA)^{-1}A^{t}A(A^tA)^{-1}A^{t} \\
	  & = & A(A^tA)^{-1}(A^{t}A)(A^tA)^{-1}A^{t} \\
	  & = & A(A^tA)^{-1}A^{t} \\
	P & = & P
	\end{array}$$
	$$\blacksquare$$
	Ahora, demostremos que $P = P^t$, esto es,
	$$\begin{array}{rcl}
	P & = & P^t \\
	  & = & (A(A^tA)^{-1}A^{t})^t \\
	  & = & ((A(A^tA)^{-1})A^{t})^t \\
	  & = & A(A(A^tA)^{-1})^t \\
	  & = & A(A^tA)^{-t}A^t \\
	  & = & A((A^tA)^t)^{-1}A^t \\
	  & = & A(A^tA)^{-1}A^t \\
	P & = & P
	\end{array}$$
	$$\blacksquare$$
\end{solucion}
\item Sea $A$ una matriz de $3\times 3$ tal que
      $$A\left(\begin{array}{r}
  1\\0\\1
\end{array}\right)= \left(\begin{array}{r}
  1\\0\\0
\end{array}\right),\quad\quad A \left(\begin{array}{r}
  1\\-2\\4
\end{array}\right)= \left(\begin{array}{r}
  1\\0\\1
\end{array}\right), \qquad  A\left(\begin{array}{r}
  -1\\1\\1
\end{array}\right)= \left(\begin{array}{r}
  0\\1\\1
\end{array}\right).$$ Calcule $A^{-1}$.
\begin{solucion}
Recordemos que para $A_{m \times n}$ y $v_1, v_2 \in \R^{n}$, siempre se cumple que
$$Av1 + Av2 = A(v_1 + v_2)$$
y recordemos también que al multiplicar dos matrices, se esta multiplicando la primera matriz con las columnas de la segunda para generar la matriz resultante, es decir, siendo $B = [v_1\ v_2\ \dots\ v_n]$,
$$AB = A[v_1\ v_2\ \dots\ v_n] = [Av_1\ Av_2\ \dots\ Av_n]$$
Procedamos con el ejercicio.\\
\\
Llamemos $v_1, v_2, v_3$ a todos los vectores que son multiplicados por $A$ en el enunciado, respectivamente y $u_1, u_2, u_3$ a sus respectivos resultados.\\
\\
Para obtener la inversa de $A$, debemos encontrar una matriz que al multiplicarla por $A$ nos de la matriz identidad.\\
\\
Busquemos entonces una forma de construir las columnas de la matriz identidad (vectores canónicos) con una combinación lineal de los vectores $u_i$. Si tuvieramos vectores más complejos, tendríamos que resolver el sistema $\alpha u_1 + \beta u_2 + \gamma u_3 = e_i$ para $i = \{1, 2, 3\}$. Sin embargo, dado que estos vectores son bastante simples, podemos hacerlo de manera analitica, esto es,
$$e_1 = u_1, \quad e_2 = , u_3-u_2+u_1, \quad e_3 = u_2 - u_1$$
Luego,
$$I = [u_1\quad u_3-u_2+u_1\quad u_2-u_1]$$
Utilizando las propiedades mencionadas al principio,
$$I = [Av_1\quad A(v_3-v_2+v_1)\quad A(v_2-v_1)]$$
$$I = A[v_1\quad v_3-v_2+v_1\quad v_2-v_1]$$
$$A^{-1} = [v_1\quad v_3-v_2+v_1\quad v_2-v_1]$$
Por lo tanto, la inversa de $A$ corresponde a
$$A^{-1} = 
\begin{bmatrix}
1 &-1 & 0 \\
0 & 3 & -2 \\
1 &-2 &3 
\end{bmatrix}$$
\end{solucion}
\item Demuestre que $T(x_1,x_2, x_3)=(x_1,x_1+x_2, x_1+x_2+x_3)$ es invertible y encuentre una f\'ormula para $T^{-1}$.
\begin{solucion}
En primer lugar determinemos la matriz de la transformación lineal $T$.\\
Sea $Ax = T(x)$,
$$A = \begin{bmatrix}
1 & 0 & 0 \\
1 & 1 & 0 \\
1 & 1 & 1
\end{bmatrix}$$
Notemos que la matriz esta escalonada "al revés". En otras palabras, $A^t$ estaría escalonada y es evidente que tiene 3 pivotes, por lo que $A$ igual tiene 3 pivotes. Al tener igual cantidad de pivotes que de columnas, la matriz $A$ es invertible.\\
\\
Ahora, busquemos $A^-1$, esto es
$$\left[
\begin{array}{ccc|ccc}
1 & 0 & 0 & 1 & 0 & 0 \\
1 & 1 & 0 & 0 & 1 & 0 \\
1 & 1 & 1 & 0 & 0 & 1
\end{array}
\right] \sim 
\left[
\begin{array}{ccc|ccc}
1 & 0 & 0 & 1 & 0 & 0 \\
0 & 1 & 0 & -1 & 1 & 0 \\
0 & 1 & 1 & -1 & 0 & 1
\end{array}
\right] \sim 
\left[
\begin{array}{ccc|ccc}
1 & 0 & 0 & 1 & 0 & 0 \\
0 & 1 & 0 & -1 & 1 & 0 \\
0 & 0 & 1 & 0 & -1 & 1
\end{array}
\right]
$$
Es decir,
$$A^{-1} = 
\begin{bmatrix}
 1 & 0 & 0 \\
-1 & 1 & 0 \\
0 & -1 & 1
\end{bmatrix}
$$
Finalmente, podemos escribir $T^{-1}(x)$ como
$$T^{-1}(x) = (x_1, -x_1 + x_2, -x_2+x_3)$$
\end{solucion}
\end{preguntas}
\end{document}